% !TEX root = ../../ctfp-print.tex

\lettrine[lhang=0.17]{圏}{が小さい(small)とは}、その対象が集合をなす場合をいう。しかし集合より大きいものが存在することは知られている。よく知られているように、すべての集合の集合は標準的な集合論（選択公理を任意に追加したツェルメロ・フレンケル集合論）の中では構成できない。したがって、すべての集合の圏は大きい(large)圏でなければならない。グロタンディーク宇宙のような数学的なトリックを使えば、集合を超えた収集を定義することができる。こうしたトリックにより、大きい圏についても議論できるようになる。

圏が\emph{局所的に小さい(locally small)}とは、任意の2つの対象間の射が集合をなす場合をいう。集合をなさない場合、いくつかの定義を再考する必要がある。特に、射を集合から選ぶことさえできないとき、射の合成とはどういう意味か？その解決策は、$\Set$の対象であるHom集合を、別の圏$\cat{V}$の\emph{対象}で置き換えることで、自身を引き上げることである。違いは、一般に対象は要素を持たないので、個々の射について語ることができなくなるという点である。\emph{豊穣圏}のすべての性質を、Hom対象全体に対して実行できる操作の観点から定義しなければならない。そのためには、Hom対象を提供する圏が追加の構造を持たなければならない――それはモノイダル圏でなければならない。このモノイダル圏を$\cat{V}$と呼ぶなら、$\cat{V}$上で豊穣化された圏$\cat{C}$について話すことができる。

大きさの問題の他にも、Hom集合を単なる集合よりも多くの構造を持つものに一般化することに興味があるかもしれない。たとえば、従来の圏には対象間の距離という概念がない。2つの対象は射で繋がっているか繋がっていないかのどちらかである。ある対象に繋がっているすべての対象はその隣人である。実生活とは異なり、圏においては友人の友人の友人も親友も同じくらい近い存在である。適切に豊穣化された圏では、対象間の距離を定義することができる。

豊穣圏の経験を積むことには、もう一つの非常に実践的な理由がある。それは、圏論の知識の非常に有用なオンラインリソースである\urlref{https://ncatlab.org/}{nLab}が、ほとんど豊穣圏の用語で書かれているからである。

\section{なぜモノイダル圏か？}

豊穣圏を構成するとき、モノイダル圏を$\Set$で置き換え、Hom対象をHom集合で置き換えたときに通常の定義を回復できるようにしなければならないことを念頭に置く必要がある。これを達成する最善の方法は、通常の定義から始めて、それらをポイントフリーな方法で――すなわち、集合の要素に名前をつけることなく――再定式化し続けることである。

合成の定義から始めよう。通常、$\cat{C}(b, c)$からの射と$\cat{C}(a, b)$からの射のペアを取り、$\cat{C}(a, c)$からの射に写す。すなわち、次のような写像である：
\[\cat{C}(b, c)\times{}\cat{C}(a, b) \to \cat{C}(a, c)\]
これは集合間の関数であり、一方は2つのHom集合のデカルト積である。この式は、デカルト積をより一般的なものに置き換えることで簡単に一般化できる。圏論的積でもよいが、さらに進んで完全に一般的なテンソル積を使うこともできる。

次は恒等射である。Hom集合から個々の要素を選ぶ代わりに、単元集合$\cat{1}$からの関数を使って定義することができる：
\[j_a \Colon \cat{1} \to \cat{C}(a, a)\]
再び、単元集合を終対象で置き換えることもできるが、さらに進んでテンソル積の単位元$i$で置き換えることができる。

見ての通り、あるモノイダル圏$\cat{V}$から取られた対象は、Hom集合の置き換えの良い候補である。

\section{モノイダル圏}

モノイダル圏については以前に述べたが、定義を再述する価値がある。モノイダル圏は、双関手であるテンソル積を定義する：
\[\otimes \Colon \cat{V}\times{}\cat{V} \to \cat{V}\]
テンソル積は結合的であることが望ましいが、自然同型を除いて結合性を満たすだけで十分である。この同型射は結合子(associator)と呼ばれる。その成分は：
\[\alpha_{a b c} \Colon (a \otimes b) \otimes c \to a \otimes (b \otimes c)\]
これは3つすべての引数において自然でなければならない。

モノイダル圏はまた、テンソル積の単位元として機能する特別な単位対象$i$を定義しなければならない。これも自然同型を除いてである。2つの同型射はそれぞれ左単位子(left unitor)と右単位子(right unitor)と呼ばれ、その成分は：
\begin{align*}
  \lambda_a & \Colon i \otimes a \to a \\
  \rho_a    & \Colon a \otimes i \to a
\end{align*}
結合子と単位子はコヒーレンス条件を満たさなければならない：

\begin{figure}[H]
  \centering
  \begin{tikzcd}[row sep=large]
    ((a \otimes b) \otimes c) \otimes d
    \arrow[d, "\alpha_{(a \otimes b)cd}"]
    \arrow[rr, "\alpha_{abc} \otimes \id_d"]
    & & (a \otimes (b \otimes c)) \otimes d
    \arrow[d, "\alpha_{a(b \otimes c)d}"] \\
    (a \otimes b) \otimes (c \otimes d)
    \arrow[rd, "\alpha_{ab(c \otimes d)}"]
    & & a \otimes ((b \otimes c) \otimes d)
    \arrow[ld, "\id_a \otimes \alpha_{bcd}"] \\
    & a \otimes (b \otimes (c \otimes d))
  \end{tikzcd}
\end{figure}

\begin{figure}[H]
  \centering
  \begin{tikzcd}[row sep=large]
    (a \otimes i) \otimes b
    \arrow[dr, "\rho_{a} \otimes \id_b"']
    \arrow[rr, "\alpha_{aib}"]
    & & a \otimes (i \otimes b)
    \arrow[dl, "\id_a \otimes \lambda_b"] \\
    & a \otimes b
  \end{tikzcd}
\end{figure}

\noindent
モノイダル圏が\newterm{対称的(symmetric)}と呼ばれるのは、次の成分を持つ自然同型が存在する場合である：
\[\gamma_{a b} \Colon a \otimes b \to b \otimes a\]
これが「二乗が1」であり：
\[\gamma_{b a} \circ \gamma_{a b} = \idarrow[a \otimes b]\]
モノイダル構造と整合性がある。

モノイダル圏の興味深い点は、テンソル積への右随伴として内部Hom（関数対象）を定義できる可能性があることである。関数対象（指数）の標準的な定義が圏論的積への右随伴を通じてであったことを思い出されるかもしれない。任意の対象のペアに対してそのような対象が存在する圏はデカルト閉(Cartesian closed)と呼ばれた。モノイダル圏における内部Homを定義する随伴は次の通りである：
\[\cat{V}(a \otimes b, c) \sim \cat{V}(a, [b, c])\]
\urlref{http://www.tac.mta.ca/tac/reprints/articles/10/tr10.pdf}{G.M.ケリー}に倣い、内部Homの記法として${[}b, c{]}$を使用する。この随伴の余単位は、その成分が評価射と呼ばれる自然変換である：
\[\varepsilon_{a b} \Colon ([a, b] \otimes a) \to b\]
テンソル積が対称でない場合、次の随伴を用いて${[}{[}a, c{]}{]}$と表す別の内部Homを定義できることに注意：
\[\cat{V}(a \otimes b, c) \sim \cat{V}(b, [[a, c]])\]
両方が定義されているモノイダル圏は\newterm{双閉(biclosed)}と呼ばれる。双閉でない圏の例は、関手合成をテンソル積とする$\Set$の自己関手の圏である。これはモナドを定義するために使った圏である。

\section{豊穣圏}

モノイダル圏$\cat{V}$上で豊穣化された圏$\cat{C}$は、Hom集合をHom対象に置き換える。$\cat{C}$内の対象$a$と$b$のすべてのペアに対して、$\cat{V}$内の対象$\cat{C}(a, b)$を対応させる。Hom対象には、Hom集合と同じ記法を使用するが、それらは射を含まないことを理解した上で使用する。一方、$\cat{V}$はHom集合と射を持つ通常の（豊穣化されていない）圏である。つまり、私たちは完全に集合から逃れたわけではなく、単にじゅうたんの下に掃き込んだだけである。

$\cat{C}$内の個々の射について語ることができないので、射の合成は$\cat{V}$内の射の族に置き換えられる：
\[\circ \Colon \cat{C}(b, c) \otimes \cat{C}(a, b) \to \cat{C}(a, c)\]

\begin{figure}[H]
  \centering
  \includegraphics[width=0.45\textwidth]{images/composition.jpg}
\end{figure}

\noindent
同様に、恒等射は$\cat{V}$内の射の族に置き換えられる：
\[j_a \Colon i \to \cat{C}(a, a)\]
ここで$i$は$\cat{V}$内のテンソル単位元である。

\begin{figure}[H]
  \centering
  \includegraphics[width=0.4\textwidth]{images/id.jpg}
\end{figure}

\noindent
合成の結合性は$\cat{V}$内の結合子の観点から定義される：

\begin{figure}[H]
  \centering
  \begin{tikzcd}[column sep=large]
    (\cat{C}(c,d) \otimes \cat{C}(b,c)) \otimes \cat{C}(a,b)
    \arrow[r, "\circ\otimes\id"]
    \arrow[dd, "\alpha"]
    & \cat{C}(b,d) \otimes \cat{C}(a,b)
    \arrow[dr, "\circ"] \\
    & & \cat{C}(a,d) \\
    \cat{C}(c,d) \otimes (\cat{C}(b,c) \otimes \cat{C}(a,b))
    \arrow[r, "\id\otimes\circ"]
    & \cat{C}(c,d) \otimes \cat{C}(a,c)
    \arrow[ur, "\circ"]
  \end{tikzcd}
\end{figure}

\noindent
単位律も同様に単位子の観点から表現される：

\begin{figure}[H]
  \centering
  \begin{subfigure}
    \centering
    \begin{tikzcd}[row sep=large]
      \cat{C}(a,b) \otimes i
      \arrow[rr, "\id \otimes j_a"]
      \arrow[dr, "\rho"]
      & & \cat{C}(a,b) \otimes \cat{C}(a,a)
      \arrow[dl, "\circ"] \\
      & \cat{C}(a,b)
    \end{tikzcd}
  \end{subfigure}
  \hspace{1cm}
  \begin{subfigure}
    \centering
    \begin{tikzcd}[row sep=large]
      i \otimes \cat{C}(a,b)
      \arrow[rr, "j_b \otimes \id"]
      \arrow[dr, "\lambda"]
      & & \cat{C}(b,b) \otimes \cat{C}(a,b)
      \arrow[dl, "\circ"] \\
      & \cat{C}(a,b)
    \end{tikzcd}
  \end{subfigure}
\end{figure}

\section{前順序}

前順序は薄い圏(thin category)として定義される。すなわち、すべてのHom集合が空か単元集合であるような圏である。空でない集合$\cat{C}(a, b)$を、$a$が$b$以下であることの証明として解釈する。そのような圏は、ただ2つの対象$0$と$1$（$\mathit{False}$と$\mathit{True}$と呼ばれることもある）を含む非常に単純なモノイダル圏上で豊穣化されたものとして解釈できる。必須の恒等射の他に、この圏には$0$から$1$に向かう1つの射があり、$0 \to 1$と呼ぶことにする。$0$と$1$の単純な算術をモデル化したテンソル積により（すなわち、唯一の非零積は$1 \otimes 1$）、単純なモノイダル構造を確立できる。この圏の単位対象は$1$である。これは厳密なモノイダル圏(strict monoidal category)、すなわち結合子と単位子が恒等射であるものである。

前順序においてHom集合は空か単元集合のどちらかであるので、それを私たちの小さな圏からのHom対象に容易に置き換えることができる。豊穣化された前順序$\cat{C}$は、任意の対象のペア$a$と$b$に対してHom対象$\cat{C}(a, b)$を持つ。$a$が$b$以下ならばこの対象は$1$であり、そうでなければ$0$である。

合成を見てみよう。任意の2つの対象のテンソル積は$0$だが、両方が$1$の場合は$1$となる。$0$の場合、合成射には2つの選択肢がある：$\idarrow[0]$か$0 \to 1$である。しかし$1$の場合、唯一の選択肢は$\idarrow[1]$である。これを関係に戻すと、$a \leqslant b$かつ$b \leqslant c$ならば$a \leqslant c$となり、これはまさに必要な推移律である。

恒等射はどうか。これは$1$から$\cat{C}(a, a)$への射である。$1$から出る射は$\idarrow[1]$のみなので、$\cat{C}(a, a)$は$1$でなければならない。これは$a \leqslant a$を意味し、前順序の反射律である。つまり、前順序を豊穣圏として実装すると、推移性と反射性の両方が自動的に強制される。

\section{距離空間}

興味深い例は\urlref{http://www.tac.mta.ca/tac/reprints/articles/1/tr1.pdf}{ウィリアム・ローヴェア}によるものである。彼は距離空間が豊穣圏を使って定義できることに気づいた。距離空間は任意の2つの対象間の距離を定義する。この距離は非負の実数である。無限大を可能な値として含めると便利である。距離が無限大の場合、出発点の対象から目標の対象へ到達する方法がない。

距離が満たさなければならない明らかな性質がいくつかある。その一つは、対象からそれ自身への距離がゼロでなければならないことである。もう一つは三角不等式：直接の距離は中間地点を経由する距離の和以下である。距離が対称であることは要求しない。これは最初は奇妙に見えるかもしれないが、ローヴェアが説明したように、一方向では上り坂を歩いていて、もう一方では下り坂を行っていると想像することができる。いずれにせよ、対称性は後で追加の制約として課すことができる。

では、距離空間はどのように圏論的な言語に翻訳できるか？Hom対象が距離であるような圏を構成しなければならない。距離は射ではなくHom対象であることに注意。Hom対象が数であることができるのは、これらの数が対象であるモノイダル圏$\cat{V}$を構成できる場合だけである。非負の実数（無限大を含む）は全順序をなすので、薄い圏として扱うことができる。2つのこのような数$x$と$y$の間の射は$x \geqslant y$のとき、かつそのときに限り存在する（注意：これは前順序の定義で伝統的に使われる方向とは逆である）。モノイダル構造は加法によって与えられ、零が単位対象として機能する。すなわち、2つの数のテンソル積はその和である。

距離空間は、このようなモノイダル圏上で豊穣化された圏である。対象$a$から$b$へのHom対象$\cat{C}(a, b)$は、$a$から$b$への距離と呼ぶ非負（場合によっては無限大）の数である。このような圏での恒等射と合成について見てみよう。

私たちの定義によれば、テンソル単位元（数ゼロ）からHom対象$\cat{C}(a, a)$への射は次の関係である：
\[0 \geqslant \cat{C}(a, a)\]
$\cat{C}(a, a)$は非負の数なので、この条件は$a$から$a$への距離が常にゼロであることを教えてくれる。確認！

合成について話そう。2つの隣接するHom対象のテンソル積$\cat{C}(b, c) \otimes \cat{C}(a, b)$から始める。テンソル積を2つの距離の和として定義した。合成は、この積から$\cat{C}(a, c)$への$\cat{V}$内の射である。$\cat{V}$内の射は、以上(greater-or-equal)の関係として定義される。言い換えると、$a$から$b$への距離と$b$から$c$への距離の和は、$a$から$c$への距離以上である。しかしこれは標準的な三角不等式に他ならない。確認！

距離空間を豊穣圏の観点で再解釈することで、三角不等式とゼロ自己距離が「タダで」得られる。

\section{豊穣関手}

関手の定義には射の写像が含まれる。豊穣化された設定では、個々の射の概念がないので、Hom対象をまとめて扱わなければならない。Hom対象はモノイダル圏$\cat{V}$内の対象であり、それらの間の射を利用できる。したがって、同じモノイダル圏$\cat{V}$上で豊穣化された圏間の豊穣関手を定義することは理にかなっている。$\cat{V}$内の射を使って、2つの豊穣圏の間のHom対象を写すことができる。

\newterm{豊穣関手(enriched functor)} $F$は、2つの圏$\cat{C}$と$\cat{D}$の間で、対象を対象に写す以外に、$\cat{C}$内の対象のすべてのペアに対して$\cat{V}$内の射を割り当てる：
\[F_{a b} \Colon \cat{C}(a, b) \to \cat{D}(F a, F b)\]
関手は構造を保存する写像である。通常の関手にとっては、合成と恒等射の保存を意味した。豊穣化された設定では、合成の保存は次の図式が可換であることを意味する：

\begin{figure}[H]
  \centering
  \begin{tikzcd}[column sep=large, row sep=large]
    \cat{C}(b,c) \otimes \cat{C}(a,b)
    \arrow[r, "\circ"]
    \arrow[d, "F_{bc} \otimes F_{ab}"]
    & \cat{C}(a,c)
    \arrow[d, "F_{ac}"] \\
    \cat{D}(F b, F c) \otimes \cat{D}(F a, F b)
    \arrow[r,  "\circ"]
    & \cat{D}(F a, F c)
  \end{tikzcd}
\end{figure}

\noindent
恒等射の保存は、恒等射を「選択する」$\cat{V}$内の射の保存に置き換えられる：

\begin{figure}[H]
  \centering
  \begin{tikzcd}[row sep=large]
    & i \arrow[dl, "j_a"'] \arrow[dr, "j_{F a}"] & \\
    \cat{C}(a,a)
    \arrow[rr, "F_{aa}"]
    & & \cat{D}(F a, F a)
  \end{tikzcd}
\end{figure}

\section{自己豊穣化}

閉じた対称モノイダル圏は、Hom集合を内部Homで置き換えることで自己豊穣化できる可能性がある（上記の定義を参照）。これを機能させるためには、内部Homの合成則を定義しなければならない。言い換えると、次のシグネチャを持つ射を実装しなければならない：
\[[b, c] \otimes [a, b] \to [a, c]\]
これは他のプログラミングタスクとあまり変わらないが、圏論ではポイントフリーな実装を使うのが通例である。まず、それが属するべき集合を特定することから始める。この場合、それはHom集合の要素である：
\[\cat{V}([b, c] \otimes [a, b], [a, c])\]
このHom集合は次と同型である：
\[\cat{V}(([b, c] \otimes [a, b]) \otimes a, c)\]
内部Hom${[}a, c{]}$を定義した随伴を使った。この新しい集合に射を構築できれば、随伴が元の集合における射を指し示し、それを合成として使用できる。この射を、利用できるいくつかの射を合成することで構築する。まず、結合子$\alpha_{{[}b, c{]}\ {[}a, b{]}\ a}$を使って左辺の式を再結合する：
\[([b, c] \otimes [a, b]) \otimes a \to [b, c] \otimes ([a, b] \otimes a)\]
続いて随伴の余単位$\varepsilon_{a b}$を使う：
\[[b, c] \otimes ([a, b] \otimes a) \to [b, c] \otimes b\]
そして余単位$\varepsilon_{b c}$を再び使って$c$に到達する。こうして射を構築した：
\[\varepsilon_{b c} \circ (\idarrow[{[b, c]}] \otimes \varepsilon_{a b}) \circ \alpha_{[b, c] [a, b] a}\]
これはHom集合の要素である：
\[\cat{V}(([b, c] \otimes [a, b]) \otimes a, c)\]
随伴は求めていた合成則を与えてくれる。

同様に、恒等射：
\[j_a \Colon i \to [a, a]\]
は次のHom集合の要素である：
\[\cat{V}(i, [a, a])\]
これは随伴を通じて次と同型である：
\[\cat{V}(i \otimes a, a)\]
このHom集合が左単位元$\lambda_a$を含むことがわかっている。随伴のもとでのその像として$j_a$を定義できる。

自己豊穣化の実践的な例はプログラミング言語における型の原型として機能する圏$\Set$である。デカルト積に関して閉じたモノイダル圏であることを以前に見た。$\Set$では、任意の2つの集合の間のHom集合はそれ自体が集合なので、$\Set$内の対象である。それが指数集合と同型であることがわかっているので、外部Homと内部Homは同値である。これで、自己豊穣化を通じて、指数集合をHom対象として使用し、合成を指数対象のデカルト積の観点で表現できることもわかった。

\section{$\cat{2}$-圏との関係}

（小）圏の圏$\Cat$の文脈で$\cat{2}$-圏について述べた。圏間の射は関手であるが、追加の構造がある：関手間の自然変換である。$\cat{2}$-圏では、対象はしばしば零セル(zero-cell)と呼ばれ、射は1セル(1-cell)、射間の射は2セル(2-cell)と呼ばれる。$\Cat$では、0セルは圏、1セルは関手、2セルは自然変換である。

しかし、2つの圏の間の関手もまた圏をなすことに注意。したがって$\Cat$では、Hom集合ではなく実際には\emph{Hom圏}を持つ。$\Set$が$\Set$上で豊穣化された圏として扱えるのと同様に、$\Cat$も$\Cat$上で豊穣化された圏として扱えることがわかる。さらに一般的に、すべての圏が$\Set$上で豊穣化されたものとして扱えるのと同様に、すべての$\cat{2}$-圏は$\Cat$上で豊穣化されたものと考えることができる。
